% TOP_PROBLEM: {"id": 1100303, "title": "Questions in Geometric Group Theory — Q 3.3", "category_id": 17, "category": {"id": 17, "name": "group_theory", "display_name": "Group Theory", "description": "Problems about groups, group actions, representations, and related algebraic structures.", "slug": "group-theory", "order_index": 17, "created_at": "2026-07-31T15:26:25.671Z"}, "status": "open", "problem_number": "AMR-010-0303", "set_id": 13, "statusQualification": "Specifies integral flat dimension; the cohomological-dimension-one theorem is a different assertion."}
% FIRST_POSTED: 2026-10-01
% ENTRY_KIND: statement_only
% UnsolvedMath ID 1100303; source catalogue code AMR-010-0303
% !TeX program = lualatex
% Definition and sources only; this file is not an LLM solution attempt.
\documentclass[11pt,a4paper]{article}
\usepackage[margin=25mm]{geometry}
\usepackage{fontspec}
\setmainfont{lmroman10-regular.otf}[BoldFont=lmroman10-bold.otf,
 ItalicFont=lmroman10-italic.otf,BoldItalicFont=lmroman10-bolditalic.otf]
\usepackage{amsmath,amssymb,mathtools}
\usepackage{unicode-math}
\setmathfont{latinmodern-math.otf}
\AtBeginDocument{\let\setminus\smallsetminus}
\usepackage{xurl,hyperref}
\hypersetup{colorlinks=true,urlcolor=blue}
\setcounter{secnumdepth}{0}
\setlength{\parindent}{0pt}
\setlength{\parskip}{5pt}
\setlength{\emergencystretch}{3em}
\begin{document}
\section*{Questions in Geometric Group Theory — Q 3.3}\label{1100303}
{\small UnsolvedMath ID 1100303 · AMR-010-0303 · Group Theory · AMR Open Problem Lists}
\subsection{Definitions and mathematical statement}
Let $G$ be a finitely generated group and let $\mathbb ZG$ be its integral group ring. Regard $\mathbb Z$ as the trivial module, with every group element acting as the identity. Its integral homological dimension is the flat dimension of this module:
\[
 \operatorname{hd}_{\mathbb Z}G
   =\min\{d:\mathbb Z\text{ has a flat }\mathbb ZG
                 \text{-resolution of length }d\}.
\]
A module is flat when tensoring with it preserves exact sequences. Equivalently, $\operatorname{hd}_{\mathbb Z}G\leq d$ means that $\operatorname{Tor}^{\mathbb ZG}_i(M,\mathbb Z)=0$ for every right $\mathbb ZG$-module $M$ and every $i>d$. These Tor groups are group homology with arbitrary module coefficients.

If $\operatorname{hd}_{\mathbb Z}G=1$, must $G$ be a free group? A free group is one having a basis subject to no relations beyond the group axioms; finite generation would make that basis finite.

The coefficient ring is the integers, not the rationals or a finite field. Homological dimension uses flat resolutions and must not be replaced by cohomological dimension, which uses projective resolutions. The group is assumed finitely generated, not finitely presented; adding finite presentation would remove the main finiteness difficulty. The equality to one excludes the dimension-zero case. The problem requests freeness of the group itself, not merely that all its finite subsets lie in free quotients.
\subsection{Short English statement}
Is every finitely generated group of integral homological dimension one free?
\subsection{Sources}
\begin{itemize}
\item[S1] ULAM.AI and the UnsolvedMath Contributors, supplied problems.json, record 1100303 (AMR-010-0303), AMR Open Problem Lists. Catalogue identity and source transcription; CC BY 4.0.
\url{https://huggingface.co/datasets/ulamai/UnsolvedMath}.
\item[S2] Bestvina - Questions in Geometric Group Theory (pdf) (2004), Question 3.3 (PDF page 10). Original source indexed by AMR-010-0303.
\url{https://www.math.utah.edu/~bestvina/eprints/questions-updated.pdf}.
\end{itemize}
\end{document}
